\section{Introduction}

The rapid growth of digital media platforms has transformed the way news is produced, consumed, and shared. Online publishers now compete not only for readership, but also for visibility across social media platforms, where article shares strongly influence audience reach and engagement. As a result, understanding the factors that drive online news popularity has become an important problem for content creators, publishers, and digital marketing teams \cite{bishop2006prml,han2011datamining}.

Predicting the popularity of an online article is a challenging task. User engagement is influenced by multiple interacting factors, including article length, language characteristics, topic, sentiment, publication timing, and media elements such as images and videos. These factors do not act independently, and the relationship between article features and popularity is often highly non-linear. This makes online news popularity an appropriate and interesting problem for machine learning and data analytics \cite{bishop2006prml,han2011datamining}.

In this project, the Online News Popularity dataset is used to study how different article characteristics relate to engagement \cite{fernandes2015dataset}. The dataset contains thousands of articles from Mashable, along with numerical features describing content structure, keyword statistics, topic probabilities, sentiment information, publication timing, and the number of social media shares. Because the dataset supports both predictive and exploratory analysis, it provides a strong basis for a broad data analytics study.

\noindent\textbf{As suggested by the professor,} the structure of this report was revised to make the analytical objectives more explicit. Instead of presenting the work mainly as a sequence of machine learning methods, the revised report frames the analysis around clearly defined tasks so that each stage of the project is linked to a specific practical question.

A key aim of this report is to move beyond simply applying algorithms and instead organize the work around meaningful analytical objectives. For that reason, the project is structured around six explicit tasks: predicting whether an article will be popular, predicting the number of shares, discovering natural article groupings, identifying engagement-related feature patterns, analyzing formatting and media usage, and investigating publication timing. This task-oriented structure makes the purpose of each analysis step clearer and aligns the technical methods with real-world questions.

The remainder of the report presents the dataset, defines the analytical tasks, explains the preparation and engineering of the data, describes the methodology used for each task, and discusses the main findings, limitations, and future directions of the study.
